%% figures/typing.tex — LaCaDiLE typing rules (Phase 1 T2)
%% Input from main.tex via %% figures/typing.tex — LaCaDiLE typing rules (Phase 1 T2)
%% Input from main.tex via %% figures/typing.tex — LaCaDiLE typing rules (Phase 1 T2)
%% Input from main.tex via %% figures/typing.tex — LaCaDiLE typing rules (Phase 1 T2)
%% Input from main.tex via \input{figures/typing.tex}
%%
%% Typing judgment: Δ; Γ ⊢ e : τ ! ε ⊣ Γ'
%%   Δ  — capability context (a set drawn from {Diff})
%%   Γ  — linear type context (input)
%%   Γ' — linear type context (output, after consumption)
%%   τ  — type
%%   ε  — effect row

\providecommand{\DiffCompat}{\mathsf{DiffCompat}}
\providecommand{\addDim}[2]{\mathsf{addDim}(#1, #2)}
\providecommand{\hastypeout}[6]{#1; #2 \vdash #3 : #4 \mathbin{!} #5 \dashv #6}

\begin{figure}[t]
\small
\begin{align*}
\addDim{d}{\tensor{\bar{d}\,}} &= \tensor{d, \bar{d}} \\
\addDim{d}{\tau_1 \to \tau_2 \mathbin{!} \varepsilon} &= \addDim{d}{\tau_1} \to \addDim{d}{\tau_2} \mathbin{!} \varepsilon \\
\addDim{d}{\tau_1 \otimes \tau_2} &= \addDim{d}{\tau_1} \otimes \addDim{d}{\tau_2} \\
\addDim{d}{\mathsf{unit}} &= \mathsf{unit} \\
\addDim{d}{\alpha} &= \alpha && \text{(type variables are unchanged)}
\end{align*}
\[
\DiffCompat \;\triangleq\; \{\mathsf{Resource}, \mathsf{Accum}\}
\qquad
\varepsilon_1 \subseteq \varepsilon_2 \;\triangleq\; \forall \mathsf{op} \in \varepsilon_1.\, \mathsf{op} \in \varepsilon_2
\]
\[
\Gamma = \Gamma_1 + \Gamma_2 \;\triangleq\;
\Gamma_1 \cap \Gamma_2 = \emptyset \;\land\; \Gamma_1 \cup \Gamma_2 = \Gamma
\quad \text{(disjoint linear split)}
\]
\[
\mathsf{ins}(\bar{d}, i, k) \;\triangleq\;
\text{the dimension list obtained from $\bar{d}$ by inserting $k$ at index $i$}
\qquad
\mathsf{rem}(\bar{d}, i) \;\triangleq\; \text{the list with index $i$ removed}
\]
\caption{Meta-functions, effect subsumption, and linear context splitting.
$\addDim{d}{\cdot}$ traverses types structurally without inspecting effect rows.
Linear context splitting is used by every multi-premise rule to thread linear
resources to exactly one sub-derivation.}
\label{fig:typing-meta}
\end{figure}

\begin{figure}[t]
\small
\begin{mathpar}
\inferrule*[right=\textsc{T-Var}]
  { }
  {\hastypeout{\Delta}{\Gamma, x{:}\tau}{x}{\tau}{\emptyset}{\Gamma}}

\inferrule*[right=\textsc{T-Unit}]
  { }
  {\hastypeout{\Delta}{\Gamma}{()}{\mathsf{unit}}{\emptyset}{\Gamma}}

\inferrule*[right=\textsc{T-Abs}]
  {\hastypeout{\Delta}{\Gamma_1, x{:}\tau_1}{e}{\tau_2}{\varepsilon}{\Gamma_2}}
  {\hastypeout{\Delta}{\Gamma_1}{\lambda x{:}\tau_1.\, e}{\tau_1 \to \tau_2 \mathbin{!} \varepsilon}{\emptyset}{\Gamma_2 \setminus \{x\}}}

\inferrule*[right=\textsc{T-App}]
  {\hastypeout{\Delta}{\Gamma_1}{e_1}{\tau_1 \to \tau_2 \mathbin{!} \varepsilon}{\varepsilon_1}{\Gamma_2} \\
   \hastypeout{\Delta}{\Gamma_2}{e_2}{\tau_1}{\varepsilon_2}{\Gamma_3}}
  {\hastypeout{\Delta}{\Gamma_1}{e_1\, e_2}{\tau_2}{\varepsilon_1 \cup \varepsilon_2 \cup \varepsilon}{\Gamma_3}}

\inferrule*[right=\textsc{T-Let}]
  {\hastypeout{\Delta}{\Gamma_1}{e_1}{\tau_1}{\varepsilon_1}{\Gamma_2} \\
   \hastypeout{\Delta}{\Gamma_2, x{:}\tau_1}{e_2}{\tau_2}{\varepsilon_2}{\Gamma_3}}
  {\hastypeout{\Delta}{\Gamma_1}{\letin{x = e_1}{e_2}}{\tau_2}{\varepsilon_1 \cup \varepsilon_2}{\Gamma_3 \setminus \{x\}}}

\inferrule*[right=\textsc{T-Copy}]
  {\hastypeout{\Delta}{\Gamma_1}{e}{\tau}{\varepsilon}{\Gamma_2}}
  {\hastypeout{\Delta}{\Gamma_1}{\mathsf{copy}(e)}{\tau \otimes \tau}{\varepsilon}{\Gamma_2}}

\inferrule*[right=\textsc{T-LetPair}]
  {\hastypeout{\Delta}{\Gamma_1}{e_1}{\tau_1 \otimes \tau_2}{\varepsilon_1}{\Gamma_2} \\
   \hastypeout{\Delta}{\Gamma_2, x{:}\tau_1, x'{:}\tau_2}{e_2}{\tau}{\varepsilon_2}{\Gamma_3}}
  {\hastypeout{\Delta}{\Gamma_1}{\letpair{x}{x'}{e_1}{e_2}}{\tau}{\varepsilon_1 \cup \varepsilon_2}{\Gamma_3 \setminus \{x, x'\}}}

\inferrule*[right=\textsc{T-Pair}]
  {\hastypeout{\Delta}{\Gamma_1}{e_1}{\tau_1}{\varepsilon_1}{\Gamma_2} \\
   \hastypeout{\Delta}{\Gamma_2}{e_2}{\tau_2}{\varepsilon_2}{\Gamma_3}}
  {\hastypeout{\Delta}{\Gamma_1}{(e_1, e_2)}{\tau_1 \otimes \tau_2}{\varepsilon_1 \cup \varepsilon_2}{\Gamma_3}}

\inferrule*[right=\textsc{T-Fst}]
  {\hastypeout{\Delta}{\Gamma_1}{e}{\tau_1 \otimes \tau_2}{\varepsilon}{\Gamma_2}}
  {\hastypeout{\Delta}{\Gamma_1}{\mathsf{fst}(e)}{\tau_1}{\varepsilon}{\Gamma_2}}

\inferrule*[right=\textsc{T-Snd}]
  {\hastypeout{\Delta}{\Gamma_1}{e}{\tau_1 \otimes \tau_2}{\varepsilon}{\Gamma_2}}
  {\hastypeout{\Delta}{\Gamma_1}{\mathsf{snd}(e)}{\tau_2}{\varepsilon}{\Gamma_2}}
\end{mathpar}
\caption{Core typing rules. Linear contexts thread left-to-right through every
multi-premise rule. $\mathsf{copy}$ is the pair-producing primitive from T0
\S1.0: its output is a linear pair whose first component is the original
(rebound) and whose second is a fresh physical copy on a new store location.
Pair destructuring via $\letpair{x}{x'}{\cdot}{\cdot}$ is the preferred
consumption form; $\mathsf{fst}$ and $\mathsf{snd}$ remain available but are
rarely used in the adjoint transformation.}
\label{fig:typing-core}
\end{figure}

\begin{figure}[t]
\small
\begin{mathpar}
\inferrule*[right=\textsc{T-Const}]
  {v \text{ is a scalar literal}}
  {\hastypeout{\Delta}{\Gamma}{\mathsf{const}(v, \bar{d})}{\tensor{\bar{d}}}{\emptyset}{\Gamma}}

\inferrule*[right=\textsc{T-Add}]
  {\hastypeout{\Delta}{\Gamma_1}{e_1}{\tensor{\bar{d}}}{\varepsilon_1}{\Gamma_2} \\
   \hastypeout{\Delta}{\Gamma_2}{e_2}{\tensor{\bar{d}}}{\varepsilon_2}{\Gamma_3}}
  {\hastypeout{\Delta}{\Gamma_1}{\mathsf{add}(e_1, e_2)}{\tensor{\bar{d}}}{\varepsilon_1 \cup \varepsilon_2}{\Gamma_3}}

\inferrule*[right=\textsc{T-Mul}]
  {\hastypeout{\Delta}{\Gamma_1}{e_1}{\tensor{\bar{d}}}{\varepsilon_1}{\Gamma_2} \\
   \hastypeout{\Delta}{\Gamma_2}{e_2}{\tensor{\bar{d}}}{\varepsilon_2}{\Gamma_3}}
  {\hastypeout{\Delta}{\Gamma_1}{\mathsf{mul}(e_1, e_2)}{\tensor{\bar{d}}}{\varepsilon_1 \cup \varepsilon_2}{\Gamma_3}}

\inferrule*[right=\textsc{T-Sum}]
  {\hastypeout{\Delta}{\Gamma_1}{e}{\tensor{\bar{d}}}{\varepsilon}{\Gamma_2} \\
   0 \leq i < |\bar{d}|}
  {\hastypeout{\Delta}{\Gamma_1}{\mathsf{sum}(e, i)}{\tensor{\mathsf{rem}(\bar{d}, i)}}{\varepsilon}{\Gamma_2}}

\inferrule*[right=\textsc{T-Expand}]
  {\hastypeout{\Delta}{\Gamma_1}{e}{\tensor{\bar{d}}}{\varepsilon}{\Gamma_2} \\
   0 \leq i \leq |\bar{d}|}
  {\hastypeout{\Delta}{\Gamma_1}{\mathsf{expand}(e, i, k)}{\tensor{\mathsf{ins}(\bar{d}, i, k)}}{\varepsilon}{\Gamma_2}}

\inferrule*[right=\textsc{T-UniformLike}]
  {\hastypeout{\Delta}{\Gamma_1}{e}{\tensor{\bar{d}}}{\varepsilon}{\Gamma_2}}
  {\hastypeout{\Delta}{\Gamma_1}{\mathsf{uniform\_like}(e, lo, hi)}{\tensor{\bar{d}}}{\varepsilon \cup \{\mathsf{Random}\}}{\Gamma_2}}
\end{mathpar}
\caption{Typing rules for the six RISC primitives. $\mathsf{add}$ and
$\mathsf{mul}$ require both operands to have equal dimension lists (enforced by
the shared $\bar{d}$ metavariable). $\mathsf{sum}$ removes the $i$-th dimension
from the type; $\mathsf{expand}$ inserts a new dimension of extent $k$ at
position $i$. $\mathsf{uniform\_like}$ is the sole stochastic primitive and is
the only way to introduce the $\mathsf{Random}$ effect.}
\label{fig:typing-prims}
\end{figure}

\begin{figure}[t]
\small
\begin{mathpar}
\inferrule*[right=\textsc{T-Perform}]
  {\hastypeout{\Delta}{\Gamma_1}{e}{\tau_{\mathsf{arg}}}{\varepsilon}{\Gamma_2} \\
   \mathsf{op} : \tau_{\mathsf{arg}} \to \tau_{\mathsf{ret}}}
  {\hastypeout{\Delta}{\Gamma_1}{\performop{op}(e)}{\tau_{\mathsf{ret}}}{\varepsilon \cup \{\mathsf{op}\}}{\Gamma_2}}

\inferrule*[right=\textsc{T-Handle}]
  {\hastypeout{\Delta}{\Gamma_1}{e}{\tau}{\varepsilon_h \cup \varepsilon_r}{\Gamma_2} \\\\
   \forall \mathsf{op}_i \in \varepsilon_h.\; \mathsf{op}_i : \tau_i^{\mathsf{arg}} \to \tau_i^{\mathsf{ret}} \\\\
   \forall i.\; \hastypeout{\Delta}{\Gamma_2, x_i{:}\tau_i^{\mathsf{arg}}, k_i{:}(\tau_i^{\mathsf{ret}} \to \tau \mathbin{!} \varepsilon_r)}{h_i}{\tau}{\varepsilon_r}{\Gamma_3}}
  {\hastypeout{\Delta}{\Gamma_1}{\handlewith{\varepsilon_h}{e}{\{\mathsf{op}_i(x_i, k_i) \to h_i\}_i}}{\tau}{\varepsilon_r}{\Gamma_3}}
\end{mathpar}
\caption{Effect system rules. $\mathsf{perform}$ introduces an effect by naming
an operation with a declared signature (e.g.\ $\mathsf{Fail}$ provides
$\mathsf{fail} : \mathsf{unit} \to \alpha$). The $\mathsf{handle}$ rule removes
the handled effects $\varepsilon_h$ from the outer effect row, leaving
$\varepsilon_r$. Each handler clause binds the operation argument $x_i$ and the
continuation $k_i$ as \emph{linear} bindings, so calling $k_i$ consumes it — the
one-shot restriction falls out of linearity. Sharing $\Gamma_3$ across all
handler clauses enforces that every branch produces the same residual context.}
\label{fig:typing-effects}
\end{figure}

\begin{figure}[t]
\small
\begin{mathpar}
\inferrule*[right=\textsc{T-Grad}]
  {\hastypeout{\Delta \cup \{\mathsf{Diff}\}}{\Gamma, x{:}\tensor{\bar{d}}}{e}{\tensor{\bar{d}'}}{\varepsilon}{\Gamma} \\\\
   \varepsilon \subseteq \DiffCompat}
  {\hastypeout{\Delta}{\Gamma}{\mathsf{grad}(\lambda x{:}\tensor{\bar{d}}.\, e)}{\tensor{\bar{d}} \to \tensor{\bar{d}'} \to \tensor{\bar{d}} \mathbin{!} \varepsilon}{\emptyset}{\Gamma}}

\inferrule*[right=\textsc{T-Vmap}]
  {\hastypeout{\Delta}{\Gamma}{f}{\tau_1 \to \tau_2 \mathbin{!} \varepsilon}{\emptyset}{\Gamma} \\
   d \text{ fresh}}
  {\hastypeout{\Delta}{\Gamma}{\mathsf{vmap}(f)}{\addDim{d}{\tau_1} \to \addDim{d}{\tau_2} \mathbin{!} \varepsilon}{\emptyset}{\Gamma}}
\end{mathpar}
\caption{AD and vectorization transforms. $\mathsf{grad}$ is defined only on
\emph{literal abstractions} $\lambda x{:}\tensor{\bar{d}}.\, e$, not on
arbitrary expressions of function type. This restriction makes the linear-use
side-condition mechanical: the premise derives $e$ with $x$ added to the
linear context and requires the output context to be $\Gamma$ unchanged, so
$x$ must be consumed exactly once by $e$. $\mathsf{grad}$ requires the
$\mathsf{Diff}$ capability in $\Delta$ and requires the body's effect row to
be a subset of $\DiffCompat = \{\mathsf{Resource}, \mathsf{Accum}\}$ —
stochastic or side-effectful bodies cannot be differentiated. The result type
is a \emph{two-argument} function: the first argument is the parameter value,
the second is the output-gradient seed (the cotangent), and the result is the
gradient with respect to the parameter. The resulting effect row is
$\varepsilon$ itself: the internal $\mathsf{Accum}$ handler introduced by the
$\mathsf{grad}$ reduction captures only $\mathsf{Accum}$, which is not present
in $\varepsilon$ to begin with (user code never emits $\mathsf{Accum}$
directly). To differentiate a function reference $f$, eta-expand first:
$\mathsf{grad}(\lambda x{:}\tensor{\bar{d}}.\, f\, x)$. $\mathsf{vmap}$ lifts
$f$ over a fresh batch dimension $d$ by applying $\addDim{d}{\cdot}$ to both
argument and result types; the effect row is preserved.}
\label{fig:typing-transforms}
\end{figure}

\noindent\textbf{Remarks on the typing rules.}
The master plan's $\&e$ borrow operator is intentionally absent. T0's tape
mechanism uses only $\mathsf{copy}$, which is sufficient for every adjoint rule
because the first component of $\mathsf{copy}$'s output pair retains the
original's ownership and the second is a fresh physical copy. A separate borrow
primitive is unnecessary for Phase 1. The $\mathsf{Fail}$ effect uses the
standard $\performop{fail}$ shape handled by $\textsc{T-Handle}$; there is no
bespoke $\textsc{T-Fail}$ rule. $\Delta$ is a non-linear set of capabilities and
threads unchanged through every rule; only $\Gamma$ is consumed. This simplifies
the substitution lemma (WS2.1) because capability-context splitting is not
required. The $\mathsf{grad}$ rule's ``$f$ uses its argument linearly'' side
condition is enforced by the linear context accounting of the premise: if $f$
is $\lambda x.\, e$ with $x$ used non-linearly in $e$, the premise derivation
fails at the inner $\textsc{T-Var}$ rule on the second use of $x$.

%%
%% Typing judgment: Δ; Γ ⊢ e : τ ! ε ⊣ Γ'
%%   Δ  — capability context (a set drawn from {Diff})
%%   Γ  — linear type context (input)
%%   Γ' — linear type context (output, after consumption)
%%   τ  — type
%%   ε  — effect row

\providecommand{\DiffCompat}{\mathsf{DiffCompat}}
\providecommand{\addDim}[2]{\mathsf{addDim}(#1, #2)}
\providecommand{\hastypeout}[6]{#1; #2 \vdash #3 : #4 \mathbin{!} #5 \dashv #6}

\begin{figure}[t]
\small
\begin{align*}
\addDim{d}{\tensor{\bar{d}\,}} &= \tensor{d, \bar{d}} \\
\addDim{d}{\tau_1 \to \tau_2 \mathbin{!} \varepsilon} &= \addDim{d}{\tau_1} \to \addDim{d}{\tau_2} \mathbin{!} \varepsilon \\
\addDim{d}{\tau_1 \otimes \tau_2} &= \addDim{d}{\tau_1} \otimes \addDim{d}{\tau_2} \\
\addDim{d}{\mathsf{unit}} &= \mathsf{unit} \\
\addDim{d}{\alpha} &= \alpha && \text{(type variables are unchanged)}
\end{align*}
\[
\DiffCompat \;\triangleq\; \{\mathsf{Resource}, \mathsf{Accum}\}
\qquad
\varepsilon_1 \subseteq \varepsilon_2 \;\triangleq\; \forall \mathsf{op} \in \varepsilon_1.\, \mathsf{op} \in \varepsilon_2
\]
\[
\Gamma = \Gamma_1 + \Gamma_2 \;\triangleq\;
\Gamma_1 \cap \Gamma_2 = \emptyset \;\land\; \Gamma_1 \cup \Gamma_2 = \Gamma
\quad \text{(disjoint linear split)}
\]
\[
\mathsf{ins}(\bar{d}, i, k) \;\triangleq\;
\text{the dimension list obtained from $\bar{d}$ by inserting $k$ at index $i$}
\qquad
\mathsf{rem}(\bar{d}, i) \;\triangleq\; \text{the list with index $i$ removed}
\]
\caption{Meta-functions, effect subsumption, and linear context splitting.
$\addDim{d}{\cdot}$ traverses types structurally without inspecting effect rows.
Linear context splitting is used by every multi-premise rule to thread linear
resources to exactly one sub-derivation.}
\label{fig:typing-meta}
\end{figure}

\begin{figure}[t]
\small
\begin{mathpar}
\inferrule*[right=\textsc{T-Var}]
  { }
  {\hastypeout{\Delta}{\Gamma, x{:}\tau}{x}{\tau}{\emptyset}{\Gamma}}

\inferrule*[right=\textsc{T-Unit}]
  { }
  {\hastypeout{\Delta}{\Gamma}{()}{\mathsf{unit}}{\emptyset}{\Gamma}}

\inferrule*[right=\textsc{T-Abs}]
  {\hastypeout{\Delta}{\Gamma_1, x{:}\tau_1}{e}{\tau_2}{\varepsilon}{\Gamma_2}}
  {\hastypeout{\Delta}{\Gamma_1}{\lambda x{:}\tau_1.\, e}{\tau_1 \to \tau_2 \mathbin{!} \varepsilon}{\emptyset}{\Gamma_2 \setminus \{x\}}}

\inferrule*[right=\textsc{T-App}]
  {\hastypeout{\Delta}{\Gamma_1}{e_1}{\tau_1 \to \tau_2 \mathbin{!} \varepsilon}{\varepsilon_1}{\Gamma_2} \\
   \hastypeout{\Delta}{\Gamma_2}{e_2}{\tau_1}{\varepsilon_2}{\Gamma_3}}
  {\hastypeout{\Delta}{\Gamma_1}{e_1\, e_2}{\tau_2}{\varepsilon_1 \cup \varepsilon_2 \cup \varepsilon}{\Gamma_3}}

\inferrule*[right=\textsc{T-Let}]
  {\hastypeout{\Delta}{\Gamma_1}{e_1}{\tau_1}{\varepsilon_1}{\Gamma_2} \\
   \hastypeout{\Delta}{\Gamma_2, x{:}\tau_1}{e_2}{\tau_2}{\varepsilon_2}{\Gamma_3}}
  {\hastypeout{\Delta}{\Gamma_1}{\letin{x = e_1}{e_2}}{\tau_2}{\varepsilon_1 \cup \varepsilon_2}{\Gamma_3 \setminus \{x\}}}

\inferrule*[right=\textsc{T-Copy}]
  {\hastypeout{\Delta}{\Gamma_1}{e}{\tau}{\varepsilon}{\Gamma_2}}
  {\hastypeout{\Delta}{\Gamma_1}{\mathsf{copy}(e)}{\tau \otimes \tau}{\varepsilon}{\Gamma_2}}

\inferrule*[right=\textsc{T-LetPair}]
  {\hastypeout{\Delta}{\Gamma_1}{e_1}{\tau_1 \otimes \tau_2}{\varepsilon_1}{\Gamma_2} \\
   \hastypeout{\Delta}{\Gamma_2, x{:}\tau_1, x'{:}\tau_2}{e_2}{\tau}{\varepsilon_2}{\Gamma_3}}
  {\hastypeout{\Delta}{\Gamma_1}{\letpair{x}{x'}{e_1}{e_2}}{\tau}{\varepsilon_1 \cup \varepsilon_2}{\Gamma_3 \setminus \{x, x'\}}}

\inferrule*[right=\textsc{T-Pair}]
  {\hastypeout{\Delta}{\Gamma_1}{e_1}{\tau_1}{\varepsilon_1}{\Gamma_2} \\
   \hastypeout{\Delta}{\Gamma_2}{e_2}{\tau_2}{\varepsilon_2}{\Gamma_3}}
  {\hastypeout{\Delta}{\Gamma_1}{(e_1, e_2)}{\tau_1 \otimes \tau_2}{\varepsilon_1 \cup \varepsilon_2}{\Gamma_3}}

\inferrule*[right=\textsc{T-Fst}]
  {\hastypeout{\Delta}{\Gamma_1}{e}{\tau_1 \otimes \tau_2}{\varepsilon}{\Gamma_2}}
  {\hastypeout{\Delta}{\Gamma_1}{\mathsf{fst}(e)}{\tau_1}{\varepsilon}{\Gamma_2}}

\inferrule*[right=\textsc{T-Snd}]
  {\hastypeout{\Delta}{\Gamma_1}{e}{\tau_1 \otimes \tau_2}{\varepsilon}{\Gamma_2}}
  {\hastypeout{\Delta}{\Gamma_1}{\mathsf{snd}(e)}{\tau_2}{\varepsilon}{\Gamma_2}}
\end{mathpar}
\caption{Core typing rules. Linear contexts thread left-to-right through every
multi-premise rule. $\mathsf{copy}$ is the pair-producing primitive from T0
\S1.0: its output is a linear pair whose first component is the original
(rebound) and whose second is a fresh physical copy on a new store location.
Pair destructuring via $\letpair{x}{x'}{\cdot}{\cdot}$ is the preferred
consumption form; $\mathsf{fst}$ and $\mathsf{snd}$ remain available but are
rarely used in the adjoint transformation.}
\label{fig:typing-core}
\end{figure}

\begin{figure}[t]
\small
\begin{mathpar}
\inferrule*[right=\textsc{T-Const}]
  {v \text{ is a scalar literal}}
  {\hastypeout{\Delta}{\Gamma}{\mathsf{const}(v, \bar{d})}{\tensor{\bar{d}}}{\emptyset}{\Gamma}}

\inferrule*[right=\textsc{T-Add}]
  {\hastypeout{\Delta}{\Gamma_1}{e_1}{\tensor{\bar{d}}}{\varepsilon_1}{\Gamma_2} \\
   \hastypeout{\Delta}{\Gamma_2}{e_2}{\tensor{\bar{d}}}{\varepsilon_2}{\Gamma_3}}
  {\hastypeout{\Delta}{\Gamma_1}{\mathsf{add}(e_1, e_2)}{\tensor{\bar{d}}}{\varepsilon_1 \cup \varepsilon_2}{\Gamma_3}}

\inferrule*[right=\textsc{T-Mul}]
  {\hastypeout{\Delta}{\Gamma_1}{e_1}{\tensor{\bar{d}}}{\varepsilon_1}{\Gamma_2} \\
   \hastypeout{\Delta}{\Gamma_2}{e_2}{\tensor{\bar{d}}}{\varepsilon_2}{\Gamma_3}}
  {\hastypeout{\Delta}{\Gamma_1}{\mathsf{mul}(e_1, e_2)}{\tensor{\bar{d}}}{\varepsilon_1 \cup \varepsilon_2}{\Gamma_3}}

\inferrule*[right=\textsc{T-Sum}]
  {\hastypeout{\Delta}{\Gamma_1}{e}{\tensor{\bar{d}}}{\varepsilon}{\Gamma_2} \\
   0 \leq i < |\bar{d}|}
  {\hastypeout{\Delta}{\Gamma_1}{\mathsf{sum}(e, i)}{\tensor{\mathsf{rem}(\bar{d}, i)}}{\varepsilon}{\Gamma_2}}

\inferrule*[right=\textsc{T-Expand}]
  {\hastypeout{\Delta}{\Gamma_1}{e}{\tensor{\bar{d}}}{\varepsilon}{\Gamma_2} \\
   0 \leq i \leq |\bar{d}|}
  {\hastypeout{\Delta}{\Gamma_1}{\mathsf{expand}(e, i, k)}{\tensor{\mathsf{ins}(\bar{d}, i, k)}}{\varepsilon}{\Gamma_2}}

\inferrule*[right=\textsc{T-UniformLike}]
  {\hastypeout{\Delta}{\Gamma_1}{e}{\tensor{\bar{d}}}{\varepsilon}{\Gamma_2}}
  {\hastypeout{\Delta}{\Gamma_1}{\mathsf{uniform\_like}(e, lo, hi)}{\tensor{\bar{d}}}{\varepsilon \cup \{\mathsf{Random}\}}{\Gamma_2}}
\end{mathpar}
\caption{Typing rules for the six RISC primitives. $\mathsf{add}$ and
$\mathsf{mul}$ require both operands to have equal dimension lists (enforced by
the shared $\bar{d}$ metavariable). $\mathsf{sum}$ removes the $i$-th dimension
from the type; $\mathsf{expand}$ inserts a new dimension of extent $k$ at
position $i$. $\mathsf{uniform\_like}$ is the sole stochastic primitive and is
the only way to introduce the $\mathsf{Random}$ effect.}
\label{fig:typing-prims}
\end{figure}

\begin{figure}[t]
\small
\begin{mathpar}
\inferrule*[right=\textsc{T-Perform}]
  {\hastypeout{\Delta}{\Gamma_1}{e}{\tau_{\mathsf{arg}}}{\varepsilon}{\Gamma_2} \\
   \mathsf{op} : \tau_{\mathsf{arg}} \to \tau_{\mathsf{ret}}}
  {\hastypeout{\Delta}{\Gamma_1}{\performop{op}(e)}{\tau_{\mathsf{ret}}}{\varepsilon \cup \{\mathsf{op}\}}{\Gamma_2}}

\inferrule*[right=\textsc{T-Handle}]
  {\hastypeout{\Delta}{\Gamma_1}{e}{\tau}{\varepsilon_h \cup \varepsilon_r}{\Gamma_2} \\\\
   \forall \mathsf{op}_i \in \varepsilon_h.\; \mathsf{op}_i : \tau_i^{\mathsf{arg}} \to \tau_i^{\mathsf{ret}} \\\\
   \forall i.\; \hastypeout{\Delta}{\Gamma_2, x_i{:}\tau_i^{\mathsf{arg}}, k_i{:}(\tau_i^{\mathsf{ret}} \to \tau \mathbin{!} \varepsilon_r)}{h_i}{\tau}{\varepsilon_r}{\Gamma_3}}
  {\hastypeout{\Delta}{\Gamma_1}{\handlewith{\varepsilon_h}{e}{\{\mathsf{op}_i(x_i, k_i) \to h_i\}_i}}{\tau}{\varepsilon_r}{\Gamma_3}}
\end{mathpar}
\caption{Effect system rules. $\mathsf{perform}$ introduces an effect by naming
an operation with a declared signature (e.g.\ $\mathsf{Fail}$ provides
$\mathsf{fail} : \mathsf{unit} \to \alpha$). The $\mathsf{handle}$ rule removes
the handled effects $\varepsilon_h$ from the outer effect row, leaving
$\varepsilon_r$. Each handler clause binds the operation argument $x_i$ and the
continuation $k_i$ as \emph{linear} bindings, so calling $k_i$ consumes it — the
one-shot restriction falls out of linearity. Sharing $\Gamma_3$ across all
handler clauses enforces that every branch produces the same residual context.}
\label{fig:typing-effects}
\end{figure}

\begin{figure}[t]
\small
\begin{mathpar}
\inferrule*[right=\textsc{T-Grad}]
  {\hastypeout{\Delta \cup \{\mathsf{Diff}\}}{\Gamma, x{:}\tensor{\bar{d}}}{e}{\tensor{\bar{d}'}}{\varepsilon}{\Gamma} \\\\
   \varepsilon \subseteq \DiffCompat}
  {\hastypeout{\Delta}{\Gamma}{\mathsf{grad}(\lambda x{:}\tensor{\bar{d}}.\, e)}{\tensor{\bar{d}} \to \tensor{\bar{d}'} \to \tensor{\bar{d}} \mathbin{!} \varepsilon}{\emptyset}{\Gamma}}

\inferrule*[right=\textsc{T-Vmap}]
  {\hastypeout{\Delta}{\Gamma}{f}{\tau_1 \to \tau_2 \mathbin{!} \varepsilon}{\emptyset}{\Gamma} \\
   d \text{ fresh}}
  {\hastypeout{\Delta}{\Gamma}{\mathsf{vmap}(f)}{\addDim{d}{\tau_1} \to \addDim{d}{\tau_2} \mathbin{!} \varepsilon}{\emptyset}{\Gamma}}
\end{mathpar}
\caption{AD and vectorization transforms. $\mathsf{grad}$ is defined only on
\emph{literal abstractions} $\lambda x{:}\tensor{\bar{d}}.\, e$, not on
arbitrary expressions of function type. This restriction makes the linear-use
side-condition mechanical: the premise derives $e$ with $x$ added to the
linear context and requires the output context to be $\Gamma$ unchanged, so
$x$ must be consumed exactly once by $e$. $\mathsf{grad}$ requires the
$\mathsf{Diff}$ capability in $\Delta$ and requires the body's effect row to
be a subset of $\DiffCompat = \{\mathsf{Resource}, \mathsf{Accum}\}$ —
stochastic or side-effectful bodies cannot be differentiated. The result type
is a \emph{two-argument} function: the first argument is the parameter value,
the second is the output-gradient seed (the cotangent), and the result is the
gradient with respect to the parameter. The resulting effect row is
$\varepsilon$ itself: the internal $\mathsf{Accum}$ handler introduced by the
$\mathsf{grad}$ reduction captures only $\mathsf{Accum}$, which is not present
in $\varepsilon$ to begin with (user code never emits $\mathsf{Accum}$
directly). To differentiate a function reference $f$, eta-expand first:
$\mathsf{grad}(\lambda x{:}\tensor{\bar{d}}.\, f\, x)$. $\mathsf{vmap}$ lifts
$f$ over a fresh batch dimension $d$ by applying $\addDim{d}{\cdot}$ to both
argument and result types; the effect row is preserved.}
\label{fig:typing-transforms}
\end{figure}

\noindent\textbf{Remarks on the typing rules.}
The master plan's $\&e$ borrow operator is intentionally absent. T0's tape
mechanism uses only $\mathsf{copy}$, which is sufficient for every adjoint rule
because the first component of $\mathsf{copy}$'s output pair retains the
original's ownership and the second is a fresh physical copy. A separate borrow
primitive is unnecessary for Phase 1. The $\mathsf{Fail}$ effect uses the
standard $\performop{fail}$ shape handled by $\textsc{T-Handle}$; there is no
bespoke $\textsc{T-Fail}$ rule. $\Delta$ is a non-linear set of capabilities and
threads unchanged through every rule; only $\Gamma$ is consumed. This simplifies
the substitution lemma (WS2.1) because capability-context splitting is not
required. The $\mathsf{grad}$ rule's ``$f$ uses its argument linearly'' side
condition is enforced by the linear context accounting of the premise: if $f$
is $\lambda x.\, e$ with $x$ used non-linearly in $e$, the premise derivation
fails at the inner $\textsc{T-Var}$ rule on the second use of $x$.

%%
%% Typing judgment: Δ; Γ ⊢ e : τ ! ε ⊣ Γ'
%%   Δ  — capability context (a set drawn from {Diff})
%%   Γ  — linear type context (input)
%%   Γ' — linear type context (output, after consumption)
%%   τ  — type
%%   ε  — effect row

\providecommand{\DiffCompat}{\mathsf{DiffCompat}}
\providecommand{\addDim}[2]{\mathsf{addDim}(#1, #2)}
\providecommand{\hastypeout}[6]{#1; #2 \vdash #3 : #4 \mathbin{!} #5 \dashv #6}

\begin{figure}[t]
\small
\begin{align*}
\addDim{d}{\tensor{\bar{d}\,}} &= \tensor{d, \bar{d}} \\
\addDim{d}{\tau_1 \to \tau_2 \mathbin{!} \varepsilon} &= \addDim{d}{\tau_1} \to \addDim{d}{\tau_2} \mathbin{!} \varepsilon \\
\addDim{d}{\tau_1 \otimes \tau_2} &= \addDim{d}{\tau_1} \otimes \addDim{d}{\tau_2} \\
\addDim{d}{\mathsf{unit}} &= \mathsf{unit} \\
\addDim{d}{\alpha} &= \alpha && \text{(type variables are unchanged)}
\end{align*}
\[
\DiffCompat \;\triangleq\; \{\mathsf{Resource}, \mathsf{Accum}\}
\qquad
\varepsilon_1 \subseteq \varepsilon_2 \;\triangleq\; \forall \mathsf{op} \in \varepsilon_1.\, \mathsf{op} \in \varepsilon_2
\]
\[
\Gamma = \Gamma_1 + \Gamma_2 \;\triangleq\;
\Gamma_1 \cap \Gamma_2 = \emptyset \;\land\; \Gamma_1 \cup \Gamma_2 = \Gamma
\quad \text{(disjoint linear split)}
\]
\[
\mathsf{ins}(\bar{d}, i, k) \;\triangleq\;
\text{the dimension list obtained from $\bar{d}$ by inserting $k$ at index $i$}
\qquad
\mathsf{rem}(\bar{d}, i) \;\triangleq\; \text{the list with index $i$ removed}
\]
\caption{Meta-functions, effect subsumption, and linear context splitting.
$\addDim{d}{\cdot}$ traverses types structurally without inspecting effect rows.
Linear context splitting is used by every multi-premise rule to thread linear
resources to exactly one sub-derivation.}
\label{fig:typing-meta}
\end{figure}

\begin{figure}[t]
\small
\begin{mathpar}
\inferrule*[right=\textsc{T-Var}]
  { }
  {\hastypeout{\Delta}{\Gamma, x{:}\tau}{x}{\tau}{\emptyset}{\Gamma}}

\inferrule*[right=\textsc{T-Unit}]
  { }
  {\hastypeout{\Delta}{\Gamma}{()}{\mathsf{unit}}{\emptyset}{\Gamma}}

\inferrule*[right=\textsc{T-Abs}]
  {\hastypeout{\Delta}{\Gamma_1, x{:}\tau_1}{e}{\tau_2}{\varepsilon}{\Gamma_2}}
  {\hastypeout{\Delta}{\Gamma_1}{\lambda x{:}\tau_1.\, e}{\tau_1 \to \tau_2 \mathbin{!} \varepsilon}{\emptyset}{\Gamma_2 \setminus \{x\}}}

\inferrule*[right=\textsc{T-App}]
  {\hastypeout{\Delta}{\Gamma_1}{e_1}{\tau_1 \to \tau_2 \mathbin{!} \varepsilon}{\varepsilon_1}{\Gamma_2} \\
   \hastypeout{\Delta}{\Gamma_2}{e_2}{\tau_1}{\varepsilon_2}{\Gamma_3}}
  {\hastypeout{\Delta}{\Gamma_1}{e_1\, e_2}{\tau_2}{\varepsilon_1 \cup \varepsilon_2 \cup \varepsilon}{\Gamma_3}}

\inferrule*[right=\textsc{T-Let}]
  {\hastypeout{\Delta}{\Gamma_1}{e_1}{\tau_1}{\varepsilon_1}{\Gamma_2} \\
   \hastypeout{\Delta}{\Gamma_2, x{:}\tau_1}{e_2}{\tau_2}{\varepsilon_2}{\Gamma_3}}
  {\hastypeout{\Delta}{\Gamma_1}{\letin{x = e_1}{e_2}}{\tau_2}{\varepsilon_1 \cup \varepsilon_2}{\Gamma_3 \setminus \{x\}}}

\inferrule*[right=\textsc{T-Copy}]
  {\hastypeout{\Delta}{\Gamma_1}{e}{\tau}{\varepsilon}{\Gamma_2}}
  {\hastypeout{\Delta}{\Gamma_1}{\mathsf{copy}(e)}{\tau \otimes \tau}{\varepsilon}{\Gamma_2}}

\inferrule*[right=\textsc{T-LetPair}]
  {\hastypeout{\Delta}{\Gamma_1}{e_1}{\tau_1 \otimes \tau_2}{\varepsilon_1}{\Gamma_2} \\
   \hastypeout{\Delta}{\Gamma_2, x{:}\tau_1, x'{:}\tau_2}{e_2}{\tau}{\varepsilon_2}{\Gamma_3}}
  {\hastypeout{\Delta}{\Gamma_1}{\letpair{x}{x'}{e_1}{e_2}}{\tau}{\varepsilon_1 \cup \varepsilon_2}{\Gamma_3 \setminus \{x, x'\}}}

\inferrule*[right=\textsc{T-Pair}]
  {\hastypeout{\Delta}{\Gamma_1}{e_1}{\tau_1}{\varepsilon_1}{\Gamma_2} \\
   \hastypeout{\Delta}{\Gamma_2}{e_2}{\tau_2}{\varepsilon_2}{\Gamma_3}}
  {\hastypeout{\Delta}{\Gamma_1}{(e_1, e_2)}{\tau_1 \otimes \tau_2}{\varepsilon_1 \cup \varepsilon_2}{\Gamma_3}}

\inferrule*[right=\textsc{T-Fst}]
  {\hastypeout{\Delta}{\Gamma_1}{e}{\tau_1 \otimes \tau_2}{\varepsilon}{\Gamma_2}}
  {\hastypeout{\Delta}{\Gamma_1}{\mathsf{fst}(e)}{\tau_1}{\varepsilon}{\Gamma_2}}

\inferrule*[right=\textsc{T-Snd}]
  {\hastypeout{\Delta}{\Gamma_1}{e}{\tau_1 \otimes \tau_2}{\varepsilon}{\Gamma_2}}
  {\hastypeout{\Delta}{\Gamma_1}{\mathsf{snd}(e)}{\tau_2}{\varepsilon}{\Gamma_2}}
\end{mathpar}
\caption{Core typing rules. Linear contexts thread left-to-right through every
multi-premise rule. $\mathsf{copy}$ is the pair-producing primitive from T0
\S1.0: its output is a linear pair whose first component is the original
(rebound) and whose second is a fresh physical copy on a new store location.
Pair destructuring via $\letpair{x}{x'}{\cdot}{\cdot}$ is the preferred
consumption form; $\mathsf{fst}$ and $\mathsf{snd}$ remain available but are
rarely used in the adjoint transformation.}
\label{fig:typing-core}
\end{figure}

\begin{figure}[t]
\small
\begin{mathpar}
\inferrule*[right=\textsc{T-Const}]
  {v \text{ is a scalar literal}}
  {\hastypeout{\Delta}{\Gamma}{\mathsf{const}(v, \bar{d})}{\tensor{\bar{d}}}{\emptyset}{\Gamma}}

\inferrule*[right=\textsc{T-Add}]
  {\hastypeout{\Delta}{\Gamma_1}{e_1}{\tensor{\bar{d}}}{\varepsilon_1}{\Gamma_2} \\
   \hastypeout{\Delta}{\Gamma_2}{e_2}{\tensor{\bar{d}}}{\varepsilon_2}{\Gamma_3}}
  {\hastypeout{\Delta}{\Gamma_1}{\mathsf{add}(e_1, e_2)}{\tensor{\bar{d}}}{\varepsilon_1 \cup \varepsilon_2}{\Gamma_3}}

\inferrule*[right=\textsc{T-Mul}]
  {\hastypeout{\Delta}{\Gamma_1}{e_1}{\tensor{\bar{d}}}{\varepsilon_1}{\Gamma_2} \\
   \hastypeout{\Delta}{\Gamma_2}{e_2}{\tensor{\bar{d}}}{\varepsilon_2}{\Gamma_3}}
  {\hastypeout{\Delta}{\Gamma_1}{\mathsf{mul}(e_1, e_2)}{\tensor{\bar{d}}}{\varepsilon_1 \cup \varepsilon_2}{\Gamma_3}}

\inferrule*[right=\textsc{T-Sum}]
  {\hastypeout{\Delta}{\Gamma_1}{e}{\tensor{\bar{d}}}{\varepsilon}{\Gamma_2} \\
   0 \leq i < |\bar{d}|}
  {\hastypeout{\Delta}{\Gamma_1}{\mathsf{sum}(e, i)}{\tensor{\mathsf{rem}(\bar{d}, i)}}{\varepsilon}{\Gamma_2}}

\inferrule*[right=\textsc{T-Expand}]
  {\hastypeout{\Delta}{\Gamma_1}{e}{\tensor{\bar{d}}}{\varepsilon}{\Gamma_2} \\
   0 \leq i \leq |\bar{d}|}
  {\hastypeout{\Delta}{\Gamma_1}{\mathsf{expand}(e, i, k)}{\tensor{\mathsf{ins}(\bar{d}, i, k)}}{\varepsilon}{\Gamma_2}}

\inferrule*[right=\textsc{T-UniformLike}]
  {\hastypeout{\Delta}{\Gamma_1}{e}{\tensor{\bar{d}}}{\varepsilon}{\Gamma_2}}
  {\hastypeout{\Delta}{\Gamma_1}{\mathsf{uniform\_like}(e, lo, hi)}{\tensor{\bar{d}}}{\varepsilon \cup \{\mathsf{Random}\}}{\Gamma_2}}
\end{mathpar}
\caption{Typing rules for the six RISC primitives. $\mathsf{add}$ and
$\mathsf{mul}$ require both operands to have equal dimension lists (enforced by
the shared $\bar{d}$ metavariable). $\mathsf{sum}$ removes the $i$-th dimension
from the type; $\mathsf{expand}$ inserts a new dimension of extent $k$ at
position $i$. $\mathsf{uniform\_like}$ is the sole stochastic primitive and is
the only way to introduce the $\mathsf{Random}$ effect.}
\label{fig:typing-prims}
\end{figure}

\begin{figure}[t]
\small
\begin{mathpar}
\inferrule*[right=\textsc{T-Perform}]
  {\hastypeout{\Delta}{\Gamma_1}{e}{\tau_{\mathsf{arg}}}{\varepsilon}{\Gamma_2} \\
   \mathsf{op} : \tau_{\mathsf{arg}} \to \tau_{\mathsf{ret}}}
  {\hastypeout{\Delta}{\Gamma_1}{\performop{op}(e)}{\tau_{\mathsf{ret}}}{\varepsilon \cup \{\mathsf{op}\}}{\Gamma_2}}

\inferrule*[right=\textsc{T-Handle}]
  {\hastypeout{\Delta}{\Gamma_1}{e}{\tau}{\varepsilon_h \cup \varepsilon_r}{\Gamma_2} \\\\
   \forall \mathsf{op}_i \in \varepsilon_h.\; \mathsf{op}_i : \tau_i^{\mathsf{arg}} \to \tau_i^{\mathsf{ret}} \\\\
   \forall i.\; \hastypeout{\Delta}{\Gamma_2, x_i{:}\tau_i^{\mathsf{arg}}, k_i{:}(\tau_i^{\mathsf{ret}} \to \tau \mathbin{!} \varepsilon_r)}{h_i}{\tau}{\varepsilon_r}{\Gamma_3}}
  {\hastypeout{\Delta}{\Gamma_1}{\handlewith{\varepsilon_h}{e}{\{\mathsf{op}_i(x_i, k_i) \to h_i\}_i}}{\tau}{\varepsilon_r}{\Gamma_3}}
\end{mathpar}
\caption{Effect system rules. $\mathsf{perform}$ introduces an effect by naming
an operation with a declared signature (e.g.\ $\mathsf{Fail}$ provides
$\mathsf{fail} : \mathsf{unit} \to \alpha$). The $\mathsf{handle}$ rule removes
the handled effects $\varepsilon_h$ from the outer effect row, leaving
$\varepsilon_r$. Each handler clause binds the operation argument $x_i$ and the
continuation $k_i$ as \emph{linear} bindings, so calling $k_i$ consumes it — the
one-shot restriction falls out of linearity. Sharing $\Gamma_3$ across all
handler clauses enforces that every branch produces the same residual context.}
\label{fig:typing-effects}
\end{figure}

\begin{figure}[t]
\small
\begin{mathpar}
\inferrule*[right=\textsc{T-Grad}]
  {\hastypeout{\Delta \cup \{\mathsf{Diff}\}}{\Gamma, x{:}\tensor{\bar{d}}}{e}{\tensor{\bar{d}'}}{\varepsilon}{\Gamma} \\\\
   \varepsilon \subseteq \DiffCompat}
  {\hastypeout{\Delta}{\Gamma}{\mathsf{grad}(\lambda x{:}\tensor{\bar{d}}.\, e)}{\tensor{\bar{d}} \to \tensor{\bar{d}'} \to \tensor{\bar{d}} \mathbin{!} \varepsilon}{\emptyset}{\Gamma}}

\inferrule*[right=\textsc{T-Vmap}]
  {\hastypeout{\Delta}{\Gamma}{f}{\tau_1 \to \tau_2 \mathbin{!} \varepsilon}{\emptyset}{\Gamma} \\
   d \text{ fresh}}
  {\hastypeout{\Delta}{\Gamma}{\mathsf{vmap}(f)}{\addDim{d}{\tau_1} \to \addDim{d}{\tau_2} \mathbin{!} \varepsilon}{\emptyset}{\Gamma}}
\end{mathpar}
\caption{AD and vectorization transforms. $\mathsf{grad}$ is defined only on
\emph{literal abstractions} $\lambda x{:}\tensor{\bar{d}}.\, e$, not on
arbitrary expressions of function type. This restriction makes the linear-use
side-condition mechanical: the premise derives $e$ with $x$ added to the
linear context and requires the output context to be $\Gamma$ unchanged, so
$x$ must be consumed exactly once by $e$. $\mathsf{grad}$ requires the
$\mathsf{Diff}$ capability in $\Delta$ and requires the body's effect row to
be a subset of $\DiffCompat = \{\mathsf{Resource}, \mathsf{Accum}\}$ —
stochastic or side-effectful bodies cannot be differentiated. The result type
is a \emph{two-argument} function: the first argument is the parameter value,
the second is the output-gradient seed (the cotangent), and the result is the
gradient with respect to the parameter. The resulting effect row is
$\varepsilon$ itself: the internal $\mathsf{Accum}$ handler introduced by the
$\mathsf{grad}$ reduction captures only $\mathsf{Accum}$, which is not present
in $\varepsilon$ to begin with (user code never emits $\mathsf{Accum}$
directly). To differentiate a function reference $f$, eta-expand first:
$\mathsf{grad}(\lambda x{:}\tensor{\bar{d}}.\, f\, x)$. $\mathsf{vmap}$ lifts
$f$ over a fresh batch dimension $d$ by applying $\addDim{d}{\cdot}$ to both
argument and result types; the effect row is preserved.}
\label{fig:typing-transforms}
\end{figure}

\noindent\textbf{Remarks on the typing rules.}
The master plan's $\&e$ borrow operator is intentionally absent. T0's tape
mechanism uses only $\mathsf{copy}$, which is sufficient for every adjoint rule
because the first component of $\mathsf{copy}$'s output pair retains the
original's ownership and the second is a fresh physical copy. A separate borrow
primitive is unnecessary for Phase 1. The $\mathsf{Fail}$ effect uses the
standard $\performop{fail}$ shape handled by $\textsc{T-Handle}$; there is no
bespoke $\textsc{T-Fail}$ rule. $\Delta$ is a non-linear set of capabilities and
threads unchanged through every rule; only $\Gamma$ is consumed. This simplifies
the substitution lemma (WS2.1) because capability-context splitting is not
required. The $\mathsf{grad}$ rule's ``$f$ uses its argument linearly'' side
condition is enforced by the linear context accounting of the premise: if $f$
is $\lambda x.\, e$ with $x$ used non-linearly in $e$, the premise derivation
fails at the inner $\textsc{T-Var}$ rule on the second use of $x$.

%%
%% Typing judgment: Δ; Γ ⊢ e : τ ! ε ⊣ Γ'
%%   Δ  — capability context (a set drawn from {Diff})
%%   Γ  — linear type context (input)
%%   Γ' — linear type context (output, after consumption)
%%   τ  — type
%%   ε  — effect row

\providecommand{\DiffCompat}{\mathsf{DiffCompat}}
\providecommand{\addDim}[2]{\mathsf{addDim}(#1, #2)}
\providecommand{\hastypeout}[6]{#1; #2 \vdash #3 : #4 \mathbin{!} #5 \dashv #6}

\begin{figure}[t]
\small
\begin{align*}
\addDim{d}{\tensor{\bar{d}\,}} &= \tensor{d, \bar{d}} \\
\addDim{d}{\tau_1 \to \tau_2 \mathbin{!} \varepsilon} &= \addDim{d}{\tau_1} \to \addDim{d}{\tau_2} \mathbin{!} \varepsilon \\
\addDim{d}{\tau_1 \otimes \tau_2} &= \addDim{d}{\tau_1} \otimes \addDim{d}{\tau_2} \\
\addDim{d}{\mathsf{unit}} &= \mathsf{unit} \\
\addDim{d}{\alpha} &= \alpha && \text{(type variables are unchanged)}
\end{align*}
\[
\DiffCompat \;\triangleq\; \{\mathsf{Resource}, \mathsf{Accum}\}
\qquad
\varepsilon_1 \subseteq \varepsilon_2 \;\triangleq\; \forall \mathsf{op} \in \varepsilon_1.\, \mathsf{op} \in \varepsilon_2
\]
\[
\Gamma = \Gamma_1 + \Gamma_2 \;\triangleq\;
\Gamma_1 \cap \Gamma_2 = \emptyset \;\land\; \Gamma_1 \cup \Gamma_2 = \Gamma
\quad \text{(disjoint linear split)}
\]
\[
\mathsf{ins}(\bar{d}, i, k) \;\triangleq\;
\text{the dimension list obtained from $\bar{d}$ by inserting $k$ at index $i$}
\qquad
\mathsf{rem}(\bar{d}, i) \;\triangleq\; \text{the list with index $i$ removed}
\]
\caption{Meta-functions, effect subsumption, and linear context splitting.
$\addDim{d}{\cdot}$ traverses types structurally without inspecting effect rows.
Linear context splitting is used by every multi-premise rule to thread linear
resources to exactly one sub-derivation.}
\label{fig:typing-meta}
\end{figure}

\begin{figure}[t]
\small
\begin{mathpar}
\inferrule*[right=\textsc{T-Var}]
  { }
  {\hastypeout{\Delta}{\Gamma, x{:}\tau}{x}{\tau}{\emptyset}{\Gamma}}

\inferrule*[right=\textsc{T-Unit}]
  { }
  {\hastypeout{\Delta}{\Gamma}{()}{\mathsf{unit}}{\emptyset}{\Gamma}}

\inferrule*[right=\textsc{T-Abs}]
  {\hastypeout{\Delta}{\Gamma_1, x{:}\tau_1}{e}{\tau_2}{\varepsilon}{\Gamma_2}}
  {\hastypeout{\Delta}{\Gamma_1}{\lambda x{:}\tau_1.\, e}{\tau_1 \to \tau_2 \mathbin{!} \varepsilon}{\emptyset}{\Gamma_2 \setminus \{x\}}}

\inferrule*[right=\textsc{T-App}]
  {\hastypeout{\Delta}{\Gamma_1}{e_1}{\tau_1 \to \tau_2 \mathbin{!} \varepsilon}{\varepsilon_1}{\Gamma_2} \\
   \hastypeout{\Delta}{\Gamma_2}{e_2}{\tau_1}{\varepsilon_2}{\Gamma_3}}
  {\hastypeout{\Delta}{\Gamma_1}{e_1\, e_2}{\tau_2}{\varepsilon_1 \cup \varepsilon_2 \cup \varepsilon}{\Gamma_3}}

\inferrule*[right=\textsc{T-Let}]
  {\hastypeout{\Delta}{\Gamma_1}{e_1}{\tau_1}{\varepsilon_1}{\Gamma_2} \\
   \hastypeout{\Delta}{\Gamma_2, x{:}\tau_1}{e_2}{\tau_2}{\varepsilon_2}{\Gamma_3}}
  {\hastypeout{\Delta}{\Gamma_1}{\letin{x = e_1}{e_2}}{\tau_2}{\varepsilon_1 \cup \varepsilon_2}{\Gamma_3 \setminus \{x\}}}

\inferrule*[right=\textsc{T-Copy}]
  {\hastypeout{\Delta}{\Gamma_1}{e}{\tau}{\varepsilon}{\Gamma_2}}
  {\hastypeout{\Delta}{\Gamma_1}{\mathsf{copy}(e)}{\tau \otimes \tau}{\varepsilon}{\Gamma_2}}

\inferrule*[right=\textsc{T-LetPair}]
  {\hastypeout{\Delta}{\Gamma_1}{e_1}{\tau_1 \otimes \tau_2}{\varepsilon_1}{\Gamma_2} \\
   \hastypeout{\Delta}{\Gamma_2, x{:}\tau_1, x'{:}\tau_2}{e_2}{\tau}{\varepsilon_2}{\Gamma_3}}
  {\hastypeout{\Delta}{\Gamma_1}{\letpair{x}{x'}{e_1}{e_2}}{\tau}{\varepsilon_1 \cup \varepsilon_2}{\Gamma_3 \setminus \{x, x'\}}}

\inferrule*[right=\textsc{T-Pair}]
  {\hastypeout{\Delta}{\Gamma_1}{e_1}{\tau_1}{\varepsilon_1}{\Gamma_2} \\
   \hastypeout{\Delta}{\Gamma_2}{e_2}{\tau_2}{\varepsilon_2}{\Gamma_3}}
  {\hastypeout{\Delta}{\Gamma_1}{(e_1, e_2)}{\tau_1 \otimes \tau_2}{\varepsilon_1 \cup \varepsilon_2}{\Gamma_3}}

\inferrule*[right=\textsc{T-Fst}]
  {\hastypeout{\Delta}{\Gamma_1}{e}{\tau_1 \otimes \tau_2}{\varepsilon}{\Gamma_2}}
  {\hastypeout{\Delta}{\Gamma_1}{\mathsf{fst}(e)}{\tau_1}{\varepsilon}{\Gamma_2}}

\inferrule*[right=\textsc{T-Snd}]
  {\hastypeout{\Delta}{\Gamma_1}{e}{\tau_1 \otimes \tau_2}{\varepsilon}{\Gamma_2}}
  {\hastypeout{\Delta}{\Gamma_1}{\mathsf{snd}(e)}{\tau_2}{\varepsilon}{\Gamma_2}}
\end{mathpar}
\caption{Core typing rules. Linear contexts thread left-to-right through every
multi-premise rule. $\mathsf{copy}$ is the pair-producing primitive from T0
\S1.0: its output is a linear pair whose first component is the original
(rebound) and whose second is a fresh physical copy on a new store location.
Pair destructuring via $\letpair{x}{x'}{\cdot}{\cdot}$ is the preferred
consumption form; $\mathsf{fst}$ and $\mathsf{snd}$ remain available but are
rarely used in the adjoint transformation.}
\label{fig:typing-core}
\end{figure}

\begin{figure}[t]
\small
\begin{mathpar}
\inferrule*[right=\textsc{T-Const}]
  {v \text{ is a scalar literal}}
  {\hastypeout{\Delta}{\Gamma}{\mathsf{const}(v, \bar{d})}{\tensor{\bar{d}}}{\emptyset}{\Gamma}}

\inferrule*[right=\textsc{T-Add}]
  {\hastypeout{\Delta}{\Gamma_1}{e_1}{\tensor{\bar{d}}}{\varepsilon_1}{\Gamma_2} \\
   \hastypeout{\Delta}{\Gamma_2}{e_2}{\tensor{\bar{d}}}{\varepsilon_2}{\Gamma_3}}
  {\hastypeout{\Delta}{\Gamma_1}{\mathsf{add}(e_1, e_2)}{\tensor{\bar{d}}}{\varepsilon_1 \cup \varepsilon_2}{\Gamma_3}}

\inferrule*[right=\textsc{T-Mul}]
  {\hastypeout{\Delta}{\Gamma_1}{e_1}{\tensor{\bar{d}}}{\varepsilon_1}{\Gamma_2} \\
   \hastypeout{\Delta}{\Gamma_2}{e_2}{\tensor{\bar{d}}}{\varepsilon_2}{\Gamma_3}}
  {\hastypeout{\Delta}{\Gamma_1}{\mathsf{mul}(e_1, e_2)}{\tensor{\bar{d}}}{\varepsilon_1 \cup \varepsilon_2}{\Gamma_3}}

\inferrule*[right=\textsc{T-Sum}]
  {\hastypeout{\Delta}{\Gamma_1}{e}{\tensor{\bar{d}}}{\varepsilon}{\Gamma_2} \\
   0 \leq i < |\bar{d}|}
  {\hastypeout{\Delta}{\Gamma_1}{\mathsf{sum}(e, i)}{\tensor{\mathsf{rem}(\bar{d}, i)}}{\varepsilon}{\Gamma_2}}

\inferrule*[right=\textsc{T-Expand}]
  {\hastypeout{\Delta}{\Gamma_1}{e}{\tensor{\bar{d}}}{\varepsilon}{\Gamma_2} \\
   0 \leq i \leq |\bar{d}|}
  {\hastypeout{\Delta}{\Gamma_1}{\mathsf{expand}(e, i, k)}{\tensor{\mathsf{ins}(\bar{d}, i, k)}}{\varepsilon}{\Gamma_2}}

\inferrule*[right=\textsc{T-UniformLike}]
  {\hastypeout{\Delta}{\Gamma_1}{e}{\tensor{\bar{d}}}{\varepsilon}{\Gamma_2}}
  {\hastypeout{\Delta}{\Gamma_1}{\mathsf{uniform\_like}(e, lo, hi)}{\tensor{\bar{d}}}{\varepsilon \cup \{\mathsf{Random}\}}{\Gamma_2}}
\end{mathpar}
\caption{Typing rules for the six RISC primitives. $\mathsf{add}$ and
$\mathsf{mul}$ require both operands to have equal dimension lists (enforced by
the shared $\bar{d}$ metavariable). $\mathsf{sum}$ removes the $i$-th dimension
from the type; $\mathsf{expand}$ inserts a new dimension of extent $k$ at
position $i$. $\mathsf{uniform\_like}$ is the sole stochastic primitive and is
the only way to introduce the $\mathsf{Random}$ effect.}
\label{fig:typing-prims}
\end{figure}

\begin{figure}[t]
\small
\begin{mathpar}
\inferrule*[right=\textsc{T-Perform}]
  {\hastypeout{\Delta}{\Gamma_1}{e}{\tau_{\mathsf{arg}}}{\varepsilon}{\Gamma_2} \\
   \mathsf{op} : \tau_{\mathsf{arg}} \to \tau_{\mathsf{ret}}}
  {\hastypeout{\Delta}{\Gamma_1}{\performop{op}(e)}{\tau_{\mathsf{ret}}}{\varepsilon \cup \{\mathsf{op}\}}{\Gamma_2}}

\inferrule*[right=\textsc{T-Handle}]
  {\hastypeout{\Delta}{\Gamma_1}{e}{\tau}{\varepsilon_h \cup \varepsilon_r}{\Gamma_2} \\\\
   \forall \mathsf{op}_i \in \varepsilon_h.\; \mathsf{op}_i : \tau_i^{\mathsf{arg}} \to \tau_i^{\mathsf{ret}} \\\\
   \forall i.\; \hastypeout{\Delta}{\Gamma_2, x_i{:}\tau_i^{\mathsf{arg}}, k_i{:}(\tau_i^{\mathsf{ret}} \to \tau \mathbin{!} \varepsilon_r)}{h_i}{\tau}{\varepsilon_r}{\Gamma_3}}
  {\hastypeout{\Delta}{\Gamma_1}{\handlewith{\varepsilon_h}{e}{\{\mathsf{op}_i(x_i, k_i) \to h_i\}_i}}{\tau}{\varepsilon_r}{\Gamma_3}}
\end{mathpar}
\caption{Effect system rules. $\mathsf{perform}$ introduces an effect by naming
an operation with a declared signature (e.g.\ $\mathsf{Fail}$ provides
$\mathsf{fail} : \mathsf{unit} \to \alpha$). The $\mathsf{handle}$ rule removes
the handled effects $\varepsilon_h$ from the outer effect row, leaving
$\varepsilon_r$. Each handler clause binds the operation argument $x_i$ and the
continuation $k_i$ as \emph{linear} bindings, so calling $k_i$ consumes it — the
one-shot restriction falls out of linearity. Sharing $\Gamma_3$ across all
handler clauses enforces that every branch produces the same residual context.}
\label{fig:typing-effects}
\end{figure}

\begin{figure}[t]
\small
\begin{mathpar}
\inferrule*[right=\textsc{T-Grad}]
  {\hastypeout{\Delta \cup \{\mathsf{Diff}\}}{\Gamma, x{:}\tensor{\bar{d}}}{e}{\tensor{\bar{d}'}}{\varepsilon}{\Gamma} \\\\
   \varepsilon \subseteq \DiffCompat}
  {\hastypeout{\Delta}{\Gamma}{\mathsf{grad}(\lambda x{:}\tensor{\bar{d}}.\, e)}{\tensor{\bar{d}} \to \tensor{\bar{d}'} \to \tensor{\bar{d}} \mathbin{!} \varepsilon}{\emptyset}{\Gamma}}

\inferrule*[right=\textsc{T-Vmap}]
  {\hastypeout{\Delta}{\Gamma}{f}{\tau_1 \to \tau_2 \mathbin{!} \varepsilon}{\emptyset}{\Gamma} \\
   d \text{ fresh}}
  {\hastypeout{\Delta}{\Gamma}{\mathsf{vmap}(f)}{\addDim{d}{\tau_1} \to \addDim{d}{\tau_2} \mathbin{!} \varepsilon}{\emptyset}{\Gamma}}
\end{mathpar}
\caption{AD and vectorization transforms. $\mathsf{grad}$ is defined only on
\emph{literal abstractions} $\lambda x{:}\tensor{\bar{d}}.\, e$, not on
arbitrary expressions of function type. This restriction makes the linear-use
side-condition mechanical: the premise derives $e$ with $x$ added to the
linear context and requires the output context to be $\Gamma$ unchanged, so
$x$ must be consumed exactly once by $e$. $\mathsf{grad}$ requires the
$\mathsf{Diff}$ capability in $\Delta$ and requires the body's effect row to
be a subset of $\DiffCompat = \{\mathsf{Resource}, \mathsf{Accum}\}$ —
stochastic or side-effectful bodies cannot be differentiated. The result type
is a \emph{two-argument} function: the first argument is the parameter value,
the second is the output-gradient seed (the cotangent), and the result is the
gradient with respect to the parameter. The resulting effect row is
$\varepsilon$ itself: the internal $\mathsf{Accum}$ handler introduced by the
$\mathsf{grad}$ reduction captures only $\mathsf{Accum}$, which is not present
in $\varepsilon$ to begin with (user code never emits $\mathsf{Accum}$
directly). To differentiate a function reference $f$, eta-expand first:
$\mathsf{grad}(\lambda x{:}\tensor{\bar{d}}.\, f\, x)$. $\mathsf{vmap}$ lifts
$f$ over a fresh batch dimension $d$ by applying $\addDim{d}{\cdot}$ to both
argument and result types; the effect row is preserved.}
\label{fig:typing-transforms}
\end{figure}

\noindent\textbf{Remarks on the typing rules.}
The master plan's $\&e$ borrow operator is intentionally absent. T0's tape
mechanism uses only $\mathsf{copy}$, which is sufficient for every adjoint rule
because the first component of $\mathsf{copy}$'s output pair retains the
original's ownership and the second is a fresh physical copy. A separate borrow
primitive is unnecessary for Phase 1. The $\mathsf{Fail}$ effect uses the
standard $\performop{fail}$ shape handled by $\textsc{T-Handle}$; there is no
bespoke $\textsc{T-Fail}$ rule. $\Delta$ is a non-linear set of capabilities and
threads unchanged through every rule; only $\Gamma$ is consumed. This simplifies
the substitution lemma (WS2.1) because capability-context splitting is not
required. The $\mathsf{grad}$ rule's ``$f$ uses its argument linearly'' side
condition is enforced by the linear context accounting of the premise: if $f$
is $\lambda x.\, e$ with $x$ used non-linearly in $e$, the premise derivation
fails at the inner $\textsc{T-Var}$ rule on the second use of $x$.
